Развитие многоядерных вычислительных систем сделало параллельные вычисление одним из способов
  ускорения задач, которые можно разделить на части. При этом производительность будет зависеть не
  только от характеристик процессора и операционной системы, но и от того, как программные потоки
  соотносятся с потоками операционной системы, как проиходит взаимодействие между ними и
  синхронизация.

В модели 1:1 каждому потоку приложения соответствует отдельный поток операционной системы. В модели
  M:N среда выполнения распределяет множество пользовательских потоков по некоторому числу потоков
  ОС. Эти модели различаются организацией планирования и использованием ресурсов.

Актуальность работы заключается в необходимости определить, как модель многопоточности влияет на
  время выполнения одной и той же вычислительной задачи и изменяется ли наблюдаемая разница на
  разных аппаратных платформах. Теоретические характеристики моделей позволяют сформулировать
  ожидаемые эффекты, однако их величина зависит от программы, размера и числа параллельных задач,
  а также конкретной вычислительной системы и требует экспериментальной проверки.

\textbf{Цель курсового проекта} — теоретически сопоставить модели многопоточности 1:1 и M:N и
  экспериментально сравнить производительность их реализаций на системах x86 и ARM при выполнении
  некоторых параллельных задач.

\textbf{Объект исследования} — процесс выполнения параллельных вычислений на современных
  вычислительных системах.

\textbf{Предмет исследования} — влияние способа отображения программных потоков на потоки
  операционной системы на производительность заданной вычислительной нагрузки.

Для достижения цели необходимо решить следующие задачи:

\begin{enumerate}
  \item рассмотреть архитектурные особенности тестовых систем x86 и ARM, существенные для
    параллельного выполнения;
  \item изучить роль операционной системы и среды выполнения в планировании программных потоков;
  \item описать модели 1:1 и M:N и сопоставить их свойства, преимущества и ограничения;
  \item выбрать вычислительную задачу и разработать сопоставимые варианты её выполнения на языке Go;
  \item проверить корректность результатов и провести повторяемые измерения на доступных системах x86 и ARM;
  \item проанализировать полученные метрики;
  \item ???? дополнительно оценить влияние операционной системы, если условия запуска в Linux и
    macOS позволят сделать содержательное сравнение.
\end{enumerate}

Экспериментальная часть будет использовать одинаковые задачи, способы их решения и входные данные
  для двух способов организации многопоточности. Для варианта M:N будут применяться обычные горутины
  Go; вариант, приближающий модель 1:1, будет строиться на закреплении горутин, выполняющих
  вычисления, за потоками ОС. Go в работе используется как средство реализации и измерения,
  а выводы будут относиться к выбранным реализациям и тестовым системам.

Основное сравнение моделей будет проводиться отдельно на компьютере с Intel Core i5-13500H с ОС
  Linux Ubuntu и на MacBook Pro с Apple M1 Pro с ОС macOS 26. Сопоставление результатов этих систем
  позволит проверить, сохраняются ли наблюдаемые тенденции на x86 и ARM. Поскольку компьютеры
  отличаются не только архитектурой набора команд, но также микроархитектурой, конфигурацией памяти
  и т.д., результаты между ними будут интерпретироваться как сравнение конкретных платформ. Linux и
  macOS рассматриваются как дополнительный фактор а не основная задача эксперимента.
